\documentclass[12pt]{article}
\usepackage[a4paper,margin=1in]{geometry}
\usepackage{amsmath,amssymb,amsfonts}
\usepackage{mathtools}
\usepackage{bm}

\title{Density-Matrix Method in QuTiP for a Single Dissipative Bosonic Mode}
\author{}
\date{}

\begin{document}
\maketitle

\section*{Purpose}
This note explains, in thesis-friendly form, what the density-matrix method does mathematically and how it is implemented in the QuTiP-based code used in this project. The discussion is specialized to the single-mode bosonic system with nonzero on-site interaction strength $U$ and nonzero single-particle loss rate $\gamma$, because that is the model actually implemented in the file
\texttt{src/bosonic\_dissipation/density\_matrix\_method.py}.

The main idea is simple:
\begin{itemize}
\item the physical state is represented by a density operator $\rho(t)$ rather than by a state vector;
\item coherent evolution is generated by the Hamiltonian;
\item dissipation is added through a Lindblad loss term;
\item after truncating the infinite bosonic Hilbert space to a finite basis, the problem becomes a finite system of coupled ordinary differential equations;
\item QuTiP solves those equations numerically and returns expectation values of observables.
\end{itemize}

\section*{Model}
For one bosonic mode with annihilation operator $\hat a$ and number operator
\[
\hat n = \hat a^\dagger \hat a,
\]
the Hamiltonian used in the code is
\[
\hat H
= \frac{U}{2}\hat a^\dagger \hat a^\dagger \hat a \hat a
= \frac{U}{2}\hat n(\hat n-1).
\]
Here $U\neq 0$ gives the on-site interaction. The open-system dynamics with single-particle loss rate $\gamma \neq 0$ is described by the Lindblad master equation
\[
\frac{d\rho}{dt}
= -i[\hat H,\rho]
+ \gamma
\left(
\hat a \rho \hat a^\dagger
- \frac{1}{2}\{\hat a^\dagger \hat a,\rho\}
\right).
\]

It is often convenient to introduce the collapse operator
\[
\hat C = \sqrt{\gamma}\,\hat a.
\]
Then the same equation is
\[
\frac{d\rho}{dt}
= -i[\hat H,\rho]
+ \hat C \rho \hat C^\dagger
- \frac{1}{2}\{\hat C^\dagger \hat C,\rho\}.
\]

\section*{Why a Density Matrix Is Used}
The density matrix is the natural object for an open quantum system because dissipation generally turns pure states into mixed states. A state vector alone cannot describe this full evolution once loss is present. The density operator
\[
\rho(t)
\]
contains both
\begin{itemize}
\item the populations $\rho_{nn}(t)$ in the number basis, and
\item the coherences $\rho_{mn}(t)$ for $m\neq n$.
\end{itemize}

For this reason, the density-matrix method evolves the full operator $\rho(t)$ rather than only a wavefunction.

\section*{Basis Truncation}
The Hilbert space of a bosonic mode is infinite, with number states
\[
|0\rangle, |1\rangle, |2\rangle, \dots.
\]
To compute numerically, the code truncates this space to
\[
|0\rangle, |1\rangle, \dots, |N-1\rangle,
\]
where $N$ is called the Hilbert-space size in the code.

Thus the density operator becomes an $N\times N$ matrix,
\[
\rho(t) \in \mathbb{C}^{N\times N}.
\]

In the present implementation:
\begin{itemize}
\item for a Fock initial state with particle number $n_0$, the minimum natural choice is $N=n_0+1$ so that $|n_0\rangle$ exists in the basis;
\item for a coherent initial state, $N$ is chosen large enough that the truncated Poisson tail is small, using the helper logic in \texttt{config.py}.
\end{itemize}

This truncation is the only approximation in the density-matrix solver itself. Once $N$ is fixed, the evolution inside that truncated space is exact for the truncated model.

\section*{Operator Matrices in the Truncated Fock Basis}
In the truncated basis $\{|0\rangle,\dots,|N-1\rangle\}$, the annihilation operator is the matrix
\[
\hat a =
\begin{pmatrix}
0 & \sqrt{1} & 0 & 0 & \cdots \\
0 & 0 & \sqrt{2} & 0 & \cdots \\
0 & 0 & 0 & \sqrt{3} & \cdots \\
\vdots & \vdots & \vdots & \ddots & \ddots \\
0 & 0 & 0 & \cdots & 0
\end{pmatrix}.
\]
Its adjoint is
\[
\hat a^\dagger =
\begin{pmatrix}
0 & 0 & 0 & \cdots & 0 \\
\sqrt{1} & 0 & 0 & \cdots & 0 \\
0 & \sqrt{2} & 0 & \cdots & 0 \\
0 & 0 & \sqrt{3} & \ddots & \vdots \\
\vdots & \vdots & \vdots & \ddots & 0
\end{pmatrix},
\]
and the number operator is diagonal,
\[
\hat n = \hat a^\dagger \hat a
= \mathrm{diag}(0,1,2,\dots,N-1).
\]

Because
\[
\hat a^\dagger \hat a^\dagger \hat a \hat a = \hat n(\hat n-1),
\]
the Hamiltonian is also diagonal in this basis:
\[
\hat H
= \frac{U}{2}\,\mathrm{diag}\!\bigl(0\cdot(-1),\,1\cdot 0,\,2\cdot 1,\dots,(N-1)(N-2)\bigr).
\]
Equivalently,
\[
H_{nn} = \frac{U}{2}\,n(n-1),
\qquad n=0,1,\dots,N-1.
\]

This immediately shows an important physical fact: the interaction term does not move population between number states. It only produces number-dependent phase rotation of the off-diagonal density-matrix elements.

\section*{What Equation Is Actually Being Solved}
Once the basis is truncated, the master equation becomes a finite-dimensional matrix differential equation:
\[
\dot{\rho}(t)
= -i\bigl(H\rho(t)-\rho(t)H\bigr)
+ \gamma\left(a\rho(t)a^\dagger - \frac{1}{2}a^\dagger a\,\rho(t) - \frac{1}{2}\rho(t)\,a^\dagger a\right).
\]
This is a linear ordinary differential equation for the $N^2$ complex entries of $\rho$.

If one writes
\[
\rho(t)=\sum_{m,n=0}^{N-1}\rho_{mn}(t)|m\rangle\langle n|,
\]
then the master equation is equivalent to a coupled system of equations for the coefficients $\rho_{mn}(t)$.

It may also be written in Liouvillian form,
\[
\frac{d\rho}{dt} = \mathcal{L}\rho,
\]
where $\mathcal{L}$ is a linear superoperator. After vectorization of the density matrix,
\[
\bm{r}(t)=\mathrm{vec}(\rho(t)),
\]
one obtains an ordinary matrix equation
\[
\frac{d\bm{r}}{dt} = L\bm{r},
\]
with
\[
L
= -i\bigl(I\otimes H - H^{T}\otimes I\bigr)
+ \gamma\left(a^{*}\otimes a - \frac12 I\otimes a^\dagger a - \frac12 (a^\dagger a)^T\otimes I\right).
\]

Conceptually, this is what QuTiP integrates internally. In other words, QuTiP is solving a linear ODE in a space of dimension $N^2$, even though the user writes the problem in operator form.

\section*{How the Initial State Is Encoded}
The code supports two initial states.

\subsection*{Fock Initial State}
For a Fock state with exactly $n_0$ particles, the code builds
\[
\rho(0)=|n_0\rangle\langle n_0|.
\]
In matrix form, this is diagonal with a single $1$ in the $(n_0,n_0)$ entry and zeros elsewhere.

\subsection*{Coherent Initial State}
For a coherent state with mean particle number $\bar n_0$, the code sets
\[
\alpha = \sqrt{\bar n_0}
\]
and constructs
\[
\rho(0)=|\alpha\rangle\langle \alpha|.
\]
In the number basis,
\[
|\alpha\rangle
= e^{-|\alpha|^2/2}\sum_{n=0}^\infty \frac{\alpha^n}{\sqrt{n!}}|n\rangle,
\]
which is then truncated to the first $N$ basis states. Therefore the coherent-state density matrix contains both diagonal populations and off-diagonal coherence terms.

\section*{How This Appears in the Actual QuTiP Code}
The key lines of the implementation are mathematically equivalent to the following steps:

\subsection*{1. Build the time grid}
\[
t_k = 0,\Delta t,2\Delta t,\dots,T.
\]
In code this is done with
\[
\texttt{time\_values = np.arange(0.0, time + dt, dt)}.
\]

\subsection*{2. Build the annihilation operator}
QuTiP creates the truncated annihilation operator with
\[
\texttt{a = qt.destroy(hilbert\_size)}.
\]
This is precisely the matrix $a$ shown above.

\subsection*{3. Build derived observables}
The code then defines
\[
\texttt{n\_op = a.dag() * a},
\]
which is $\hat n$, and
\[
\texttt{factorial\_second\_moment\_op = (a.dag() ** 2) * (a ** 2)},
\]
which is
\[
\hat a^{\dagger 2}\hat a^2 = \hat n(\hat n-1).
\]

\subsection*{4. Build the Hamiltonian}
The implemented Hamiltonian is
\[
\texttt{hamiltonian = 0.5 * interaction\_strength * (a.dag() ** 2) * (a ** 2)}.
\]
Therefore, with \texttt{interaction\_strength = U},
\[
\hat H = \frac{U}{2}\hat a^\dagger \hat a^\dagger \hat a \hat a.
\]

\subsection*{5. Build the collapse operators}
The loss channel is entered as
\[
\texttt{collapse\_operators = [np.sqrt(gamma) * a]}.
\]
Thus QuTiP receives the collapse operator
\[
\hat C = \sqrt{\gamma}\,\hat a,
\]
which generates the Lindblad dissipator for single-particle loss.

\subsection*{6. Build the initial density matrix}
The helper routine constructs either
\[
\rho_0 = |n_0\rangle\langle n_0|
\]
or
\[
\rho_0 = |\alpha\rangle\langle \alpha|,
\qquad \alpha=\sqrt{\bar n_0}.
\]

\subsection*{7. Integrate the master equation}
The core QuTiP call is
\[
\texttt{qt.mesolve(hamiltonian, rho0, time\_values, collapse\_operators, [a, n\_op, factorial\_second\_moment\_op])}.
\]
Mathematically, \texttt{mesolve} evolves $\rho(t)$ according to the Lindblad master equation and simultaneously computes
\[
\langle \hat a \rangle(t), \qquad
\langle \hat n \rangle(t), \qquad
\langle \hat a^{\dagger 2}\hat a^2 \rangle(t)
\]
at every time point in the grid.

\section*{What QuTiP Returns}
At each time $t_k$, QuTiP returns expectation values of the requested operators:
\[
\langle \hat O \rangle(t_k) = \mathrm{Tr}\!\bigl(\rho(t_k)\hat O\bigr).
\]
In this project the returned quantities are:
\begin{itemize}
\item \textbf{Field amplitude}
\[
g_1(t) \equiv \langle \hat a \rangle(t).
\]
Strictly speaking, this object is the first moment of the field. The code stores it in the variable \texttt{g1}.
\item \textbf{Mean particle number}
\[
\langle \hat n \rangle(t).
\]
\item \textbf{Factorial second moment}
\[
F_2(t)=\langle \hat a^{\dagger 2}\hat a^2\rangle(t)=\langle \hat n(\hat n-1)\rangle(t).
\]
\end{itemize}

From these, the code computes
\[
\mathrm{Var}(\hat n)
= \langle \hat n^2\rangle - \langle \hat n\rangle^2.
\]
Since
\[
\hat n^2 = \hat n(\hat n-1)+\hat n,
\]
one has
\[
\langle \hat n^2\rangle = F_2 + \langle \hat n\rangle.
\]
Therefore the variance is reconstructed as
\[
\mathrm{Var}(\hat n)
= F_2 + \langle \hat n\rangle - \langle \hat n\rangle^2,
\]
which is exactly what the code evaluates.

The equal-time second-order coherence is then computed as
\[
g^{(2)}(t)
= \frac{\langle \hat a^{\dagger 2}\hat a^2\rangle(t)}{\langle \hat n\rangle(t)^2}
= \frac{F_2(t)}{\langle \hat n\rangle(t)^2},
\]
whenever the denominator is nonzero.

\section*{What the Interaction $U$ Does Mathematically}
Because
\[
\hat H = \frac{U}{2}\hat n(\hat n-1)
\]
is diagonal in the number basis, it does not directly change the diagonal populations by itself. Instead, it generates phase evolution of the off-diagonal terms:
\[
\rho_{mn}(t)
\sim
e^{-i(E_m-E_n)t}\rho_{mn}(0),
\qquad
E_n = \frac{U}{2}n(n-1),
\]
before dissipation is included.

So the role of $U$ is:
\begin{itemize}
\item it changes relative phases between different number sectors;
\item it affects coherence-sensitive observables, such as $\langle \hat a\rangle$;
\item it does not directly reshuffle occupation probabilities through the Hamiltonian part alone.
\end{itemize}

The role of $\gamma$ is different:
\begin{itemize}
\item it removes population through the loss channel;
\item it damps both populations and coherences;
\item it makes the evolution non-unitary and generally mixed-state.
\end{itemize}

\section*{A Small Matrix Example}
To make the procedure concrete, suppose the Hilbert space is truncated at $N=4$, so the basis is
\[
|0\rangle, |1\rangle, |2\rangle, |3\rangle.
\]
Then
\[
a =
\begin{pmatrix}
0 & 1 & 0 & 0 \\
0 & 0 & \sqrt{2} & 0 \\
0 & 0 & 0 & \sqrt{3} \\
0 & 0 & 0 & 0
\end{pmatrix},
\qquad
n =
\begin{pmatrix}
0 & 0 & 0 & 0 \\
0 & 1 & 0 & 0 \\
0 & 0 & 2 & 0 \\
0 & 0 & 0 & 3
\end{pmatrix}.
\]
The Hamiltonian is
\[
H = \frac{U}{2}n(n-1)
= \mathrm{diag}(0,0,U,3U).
\]
The collapse operator is
\[
C=\sqrt{\gamma}\,a.
\]

The density matrix is
\[
\rho(t)=
\begin{pmatrix}
\rho_{00} & \rho_{01} & \rho_{02} & \rho_{03} \\
\rho_{10} & \rho_{11} & \rho_{12} & \rho_{13} \\
\rho_{20} & \rho_{21} & \rho_{22} & \rho_{23} \\
\rho_{30} & \rho_{31} & \rho_{32} & \rho_{33}
\end{pmatrix},
\]
and the master equation becomes a coupled system for these 16 entries. QuTiP constructs and solves this system numerically.

\section*{Interpretation of the Numerical Method}
The density-matrix method used here is therefore not a stochastic method and not a sampling method. It does not average over trajectories. Instead, it directly solves the deterministic master equation for the full reduced density operator in the truncated basis.

Its main advantages are:
\begin{itemize}
\item it is conceptually direct;
\item it treats mixed states and dissipation exactly within the chosen truncation;
\item it gives expectation values without statistical noise.
\end{itemize}

Its main computational cost is that the unknown object is an $N\times N$ matrix, so the state dimension scales like $N^2$ rather than $N$. This is why the method becomes expensive when the required Hilbert-space truncation is large.

\section*{Summary}
For the single-mode bosonic model with nonzero $U$ and nonzero $\gamma$, the density-matrix method implemented in QuTiP proceeds as follows:
\begin{enumerate}
\item choose a finite Fock basis of size $N$;
\item construct the annihilation matrix $a$ and number matrix $n=a^\dagger a$;
\item build the Hamiltonian
\[
H=\frac{U}{2}a^\dagger a^\dagger a a;
\]
\item build the collapse operator
\[
C=\sqrt{\gamma}\,a;
\]
\item construct the initial density matrix $\rho_0$;
\item solve the Lindblad master equation
\[
\dot{\rho}=-i[H,\rho]+C\rho C^\dagger-\frac12\{C^\dagger C,\rho\};
\]
\item compute observables from
\[
\langle O\rangle(t)=\mathrm{Tr}[\rho(t)O].
\]
\end{enumerate}

Mathematically, the method turns the open quantum problem into a finite linear ODE for the entries of $\rho(t)$, or equivalently for the vectorized density matrix. In the code, QuTiP handles this operator-to-matrix machinery internally, while the user specifies the physical ingredients directly in terms of $H$, $\rho_0$, and the collapse operators.

\end{document}
